\documentclass[tikz,border=8pt]{standalone}
\usepackage{tikz}
\usepackage{amsmath}
\usepackage{amssymb}
\usepackage{pgfplots}
\pgfplotsset{compat=1.18}
\usepackage{ctex}
\usetikzlibrary{calc,positioning,arrows.meta,matrix,trees}

% LC 136 Single Number XOR
% 数组 [4,1,2,1,2]，累计 XOR：4→5→7→6→4，最终结果 4

\begin{document}
\begin{tikzpicture}[
  every node/.style={font=\small},
  cell/.style={draw=gray!60, rectangle, minimum width=0.75cm, minimum height=0.6cm, thick, fill=white},
  bitcell/.style={draw=gray!50, rectangle, minimum width=0.45cm, minimum height=0.5cm, thin, fill=white},
  curbit/.style={draw=blue!60, rectangle, minimum width=0.45cm, minimum height=0.5cm, thin, fill=blue!12},
  resbit/.style={draw=green!60!black, rectangle, minimum width=0.45cm, minimum height=0.5cm, thin, fill=green!12},
  xorbox/.style={draw=orange!70!black, rectangle, minimum width=0.75cm, minimum height=0.6cm, thick, fill=orange!12},
  ptr/.style={->,>=Stealth,thick},
]

%% 标题
\node[font=\large\bfseries] at (5.5, 1.5)
  {XOR 找单数（LC 136 Single Number）};
\node[font=\small,gray] at (5.5, 0.9)
  {数组 [4,1,2,1,2]，逐步 XOR，重复元素两两抵消，剩余即答案};

%% 输入数组
\node[font=\scriptsize,gray,anchor=east] at (0.8, 0.2) {nums:};
\foreach \i/\v in {0/4, 1/1, 2/2, 3/1, 4/2} {
  \node[cell,font=\ttfamily] at (1.2+\i*0.85, 0.2) {\v};
  \node[font=\tiny,gray] at (1.2+\i*0.85, -0.22) {\i};
}

%% XOR 步骤表（4步累计）
\node[font=\small\bfseries] at (5.5, -0.8) {逐步 XOR 累计};

% 表头
\node[font=\scriptsize\bfseries,gray] at (1.8, -1.35) {步骤};
\node[font=\scriptsize\bfseries,gray] at (3.5, -1.35) {操作};
\node[font=\scriptsize\bfseries,gray] at (6.5, -1.35) {二进制（4位）};
\node[font=\scriptsize\bfseries,gray] at (9.5, -1.35) {结果};

\draw[gray!40] (0.8, -1.55) -- (11.0, -1.55);

% 步骤0：初始 result=0
\node[font=\scriptsize] at (1.8, -1.9) {初始};
\node[font=\scriptsize,gray] at (3.5, -1.9) {result = 0};
% 二进制
\foreach \i/\b in {0/0,1/0,2/0,3/0} {
  \node[bitcell] at (5.5+\i*0.5, -1.9) {\b};
}
\node[font=\scriptsize,gray] at (9.5, -1.9) {0};

% 步骤1：XOR 4 = 0100
\node[font=\scriptsize] at (1.8, -2.6) {i=0};
\node[font=\scriptsize] at (3.5, -2.6) {0 XOR \textbf{4}};
\foreach \i/\b in {0/0,1/1,2/0,3/0} {
  \node[curbit] at (5.5+\i*0.5, -2.6) {\b};
}
\node[font=\scriptsize,xorbox] at (9.5, -2.6) {4};
\node[font=\tiny,gray] at (8.3, -2.6) {0100};

% 步骤2：XOR 1 = 0001 → result=5
\node[font=\scriptsize] at (1.8, -3.3) {i=1};
\node[font=\scriptsize] at (3.5, -3.3) {4 XOR \textbf{1}};
\foreach \i/\b in {0/0,1/1,2/0,3/1} {
  \node[curbit] at (5.5+\i*0.5, -3.3) {\b};
}
\node[font=\scriptsize,xorbox] at (9.5, -3.3) {5};
\node[font=\tiny,gray] at (8.3, -3.3) {0101};

% 步骤3：XOR 2 = 0010 → result=7
\node[font=\scriptsize] at (1.8, -4.0) {i=2};
\node[font=\scriptsize] at (3.5, -4.0) {5 XOR \textbf{2}};
\foreach \i/\b in {0/0,1/1,2/1,3/1} {
  \node[curbit] at (5.5+\i*0.5, -4.0) {\b};
}
\node[font=\scriptsize,xorbox] at (9.5, -4.0) {7};
\node[font=\tiny,gray] at (8.3, -4.0) {0111};

% 步骤4：XOR 1 = 0001 → result=6
\node[font=\scriptsize] at (1.8, -4.7) {i=3};
\node[font=\scriptsize] at (3.5, -4.7) {7 XOR \textbf{1}};
\foreach \i/\b in {0/0,1/1,2/1,3/0} {
  \node[curbit] at (5.5+\i*0.5, -4.7) {\b};
}
\node[font=\scriptsize,xorbox] at (9.5, -4.7) {6};
\node[font=\tiny,gray] at (8.3, -4.7) {0110};

% 步骤5：XOR 2 = 0010 → result=4
\node[font=\scriptsize] at (1.8, -5.4) {i=4};
\node[font=\scriptsize] at (3.5, -5.4) {6 XOR \textbf{2}};
\foreach \i/\b in {0/0,1/1,2/0,3/0} {
  \node[resbit] at (5.5+\i*0.5, -5.4) {\b};
}
\node[font=\scriptsize,draw=green!60!black,fill=green!15,rounded corners=2pt,
      inner sep=2pt] at (9.5, -5.4) {\textbf{4}};
\node[font=\tiny,gray] at (8.3, -5.4) {0100};

\draw[gray!40] (0.8, -5.7) -- (11.0, -5.7);

%% XOR 真值表
\node[font=\small\bfseries] at (2.0, -6.3) {XOR 真值表};

\node[font=\scriptsize,gray] at (1.0, -6.8) {A};
\node[font=\scriptsize,gray] at (1.6, -6.8) {B};
\node[font=\scriptsize,gray] at (2.3, -6.8) {A$\oplus$B};
\draw[gray!40] (0.65, -7.0) -- (2.7, -7.0);
\foreach \a/\b/\r in {0/0/0, 0/1/1, 1/0/1, 1/1/0} {
  \pgfmathsetmacro{\rr}{-7.3 - (\a*2+\b)*0.4}
  \node[font=\scriptsize] at (1.0, \rr) {\a};
  \node[font=\scriptsize] at (1.6, \rr) {\b};
  \node[font=\scriptsize] at (2.3, \rr) {\r};
}

\node[draw=gray!40,fill=white,thin,rounded corners=2pt,
      font=\scriptsize,inner sep=5pt,align=left]
  at (7.5, -6.8)
  {关键性质：\\[3pt]
   $a \oplus a = 0$（相同数异或得0）\\[2pt]
   $a \oplus 0 = a$（与0异或不变）\\[2pt]
   XOR 满足交换律和结合律};

%% 结果框
\node[draw=green!60!black,fill=green!8,very thick,rounded corners=4pt,
      font=\large\bfseries]
  at (5.5, -9.0) {答案 = 4（唯一出现一次的数）};

%% 算法说明
\node[draw=gray!50,fill=white,rounded corners=3pt,font=\small,
      inner sep=8pt,align=left]
  at (5.5, -10.5)
  {\textbf{Single Number O(n) 时间，O(1) 空间：}\\[4pt]
   \texttt{result = 0}，遍历每个元素 x：\texttt{result \^{}= x}\\[2pt]
   成对出现的数异或后归零，只剩单次出现的数};

\end{tikzpicture}
\end{document}
